\documentclass[12pt]{article}
\usepackage[T1]{fontenc}
\usepackage[utf8]{inputenc}
\usepackage{lmodern}
\usepackage{textcomp}
\usepackage{enumitem}
\usepackage{geometry}
\usepackage{graphicx}
\usepackage{amsmath, amssymb}
\geometry{margin=1in}
\begin{document}
\begin{center}
Astronomy - Graduate Level Assessment 22 \
Total Marks: 40 \
Answer all questions.
\end{center}

\section*{Multiple Choice (10 marks)}
\begin{enumerate}[label=\arabic*.]
\item The lithosphere of a planet is the layer that consists of
\begin{enumerate}[label=(\alph*)]
    \item the softer rocky material of the mantle.
    \item the lava that comes out of volcanoes.
    \item material between the crust and the mantle.
    \item the rigid rocky material of the crust and uppermost portion of the mantle.
\end{enumerate}
\item How do we know how old the Earth is?
\begin{enumerate}[label=(\alph*)]
    \item From the layering of materials within the Earth.
    \item From fossils of ancient life.
    \item From the cratering history of Earth's surface.
    \item From radioactive dating of rocks and meteorites.
\end{enumerate}
\item If the Moon is setting at noon the phase of the Moon must be
\begin{enumerate}[label=(\alph*)]
    \item third quarter.
    \item waning crescent.
    \item waxing crescent.
    \item full.
\end{enumerate}
\item You are pushing a truck along a road. Would it be easier to accelerate this truck on Mars? Why? (Assume there is no friction)
\begin{enumerate}[label=(\alph*)]
    \item It would be harder since the truck is heavier on Mars.
    \item It would be easier since the truck is lighter on Mars.
    \item It would be harder since the truck is lighter on Mars.
    \item It would be the same no matter where you are.
\end{enumerate}
\item Which of the following lists the ingredients of the solar nebula from highest to lowest percentage of mass in the nebula?
\begin{enumerate}[label=(\alph*)]
    \item hydrogen compounds (H$_{2}$OCH$_{4}$NH$_{3}$) light gases (H He) metals rocks
    \item hydrogen compounds (H$_{2}$OCH$_{4}$NH$_{3}$) light gases (H He) rocks metals
    \item light gases (H He) hydrogen compounds (H$_{2}$OCH$_{4}$NH$_{3}$) metals rocks
    \item light gases (H He) hydrogen compounds (H$_{2}$OCH$_{4}$NH$_{3}$) rocks metals
\end{enumerate}
\end{enumerate}

\section*{True / False (10 marks)}
\textit{Answer True or False}
\begin{enumerate}[label=\arabic*.]
\setcounter{enumi}{5}
\item True or False: Planetary rings can form from the dismantling of small moons due to impacts from larger celestial bodies.
\item True or False: Venus shows no signs of volcanic activity, indicating a geologically inactive surface.
\item True or False: Barnard's Star is the second closest star system to Earth after Proxima Centauri.
\item True or False: The Sun's corona cannot be seen because it emits only infrared radiation, which is not visible to the naked eye.
\item True or False: In$_{20}$,000 years, the Moon will be closer to the Earth and the length of the Earth's day will decrease due to gravitational effects.
\end{enumerate}

\section*{Long Form Response (20 marks)}
\textit{Answer each question with a clear and complete explanation.}
\begin{enumerate}[label=\arabic*.]
\setcounter{enumi}{10}
\item Discuss the significance of Betelgeuse within the constellation Orion, including its characteristics, role in stellar evolution, and its cultural impact in astronomy.
\item Discuss the implications of Kepler's laws of planetary motion and evaluate why the inverse square law of gravitational attraction is not a direct consequence of these laws.
\end{enumerate}

\vfill
\noindent\textit{End of Paper}
\end{document}